\documentclass[tikz,border=10pt]{standalone}
\input{../../../docs/tikz/dag_preamble.tex}

\begin{document}

\begin{tikzpicture}[
  x=1cm,
  y=1cm,
  every node/.append style={outer sep=4pt, font=\sffamily\normalsize}
]

\dagPaperMetadata{Iterative Local Voting for Collective Decision-making in Continuous Spaces}{N. Garg, V. Kamble, A. Goel, D. Marn, K. Munagala}{JAIR 64:315--355}{2019}{https://doi.org/10.1613/jair.1.11358}

\dagPaperLegendRightOfMetadata{
\node[dag_result, dag_template_legend] (legRes) at (0,0) {formalized\\result};
\node[dag_lemma, dag_template_legend] (legLem) at (4.2,0) {formalized\\lemma};
\node[dag_model, dag_template_legend] (legDef) at (8.4,0) {formalized\\definition};
\node[dag_caveat_legend] (legCav) at (12.6,0) {formalized\\with caveat};
\node[dag_partial, dag_template_legend] (legPart) at (0,-2.6) {partially\\formalized};
\node[dag_scaffold, dag_template_legend] (legScaf) at (4.2,-2.6) {scaffold};
\node[dag_unformalized, dag_template_legend] (legNot) at (8.4,-2.6) {not started /\\not formalized};
}
\daglegend{(legRes)(legLem)(legDef)(legCav)(legPart)(legScaf)(legNot)}{Legend}

\begin{dagPaperBody}

\node[dag_model] (C123) at (0,-4.4) {
  \textbf{C1--C3 and Algorithm 1} \\
  Convex bounded space, unique ideals, bounded density, radius $r_0/t$
};
\node[dag_model] (Definitions) at (8.4,-4.4) {
  \textbf{Definitions 1--3} \\
  $L^p$, weighted Euclidean, and decomposable utility classes
};
\node[dag_lemma] (Appendix) at (16.8,-4.4) {
  \textbf{Appendix analytic support} \\
  Lemmas 1--4, including the corrected block-event estimate
};

\node[dag_result] (Theorem1) at (0,-10.2) {
  \textbf{Theorem 1} \\
  Six finite-coordinate iid norm/model cases
};
\node[dag_result] (Theorem2) at (8.4,-10.2) {
  \textbf{Theorem 2} \\
  Finite-exponent Model B Holder-dual cases
};
\node[dag_result, text width=5.3cm] (Proposition1) at (16.8,-10.2) {
  \textbf{Proposition 1} \\
  Concrete Definition 2 joint-law Model A and B convergence;\\
  Model A uses canonical water-filling
};

\node[dag_result, text width=5.8cm] (Proposition2) at (4.2,-16.4) {
  \textbf{Proposition 2} \\
  Coordinatewise median convergence;\\
  Model B uses full-radius coordinate steps
};
\node[dag_result, text width=5.3cm] (Theorem3Constrained) at (12.6,-16.4) {
  \textbf{Theorem 3} \\
  Constrained-space alternative
};
\node[dag_result, text width=4.6cm] (Theorem3Full) at (18.8,-16.4) {
  \textbf{Theorem 3} \\
  Printed endpoint in full space
};

\node[dag_result, text width=7.6cm, minimum width=8.0cm] (PaperStatus) at (9.4,-25.2) {
  \textbf{Paper status} \\
  Formalized: all selected results are proved; Proposition 2 and Theorem 3
  use the source clarifications recorded in the validation report.
};

\draw[dag_arrow] (C123) -- (Theorem1);
\draw[dag_arrow] (C123) -- (Theorem2);
\draw[dag_arrow] (C123.east) -- (Proposition1.west);
\draw[dag_arrow] (Definitions.south west) -- (Theorem1.north east);
\draw[dag_arrow] (Definitions) -- (Theorem2);
\draw[dag_arrow] (Definitions.south east) -- (Proposition1.north west);
\draw[dag_arrow] (Definitions.south) -- (Proposition2.north east);
\draw[dag_arrow] (Appendix.south west) -- (Theorem2.north east);
\draw[dag_arrow] (Appendix) -- (Proposition1);
\draw[dag_arrow] (C123.south east) -- (Theorem3Constrained.north west);
\draw[dag_arrow] (Definitions.south east) -- (Theorem3Constrained.north west);
\draw[dag_arrow] (Theorem3Constrained) -- (Theorem3Full);
\draw[dag_dashed_arrow] (Theorem1) -- (PaperStatus.north west);
\draw[dag_dashed_arrow] (Theorem2) -- (PaperStatus.north);
\draw[dag_dashed_arrow] (Proposition1) -- (PaperStatus.north east);
\draw[dag_dashed_arrow] (Proposition2) -- (PaperStatus.west);
\draw[dag_dashed_arrow] (Theorem3Constrained.south) -- (PaperStatus.north east);
\draw[dag_dashed_arrow] (Theorem3Full.south west) -- (PaperStatus.east);
\end{dagPaperBody}
\end{tikzpicture}
\end{document}
